% ==============================================================================
% CHAPTER 4: RESULT AND DISCUSSION
% ==============================================================================

\chapter{Result and Discussion}

This chapter presents the empirical results of the deployed Integrated agricultural product data System with particular emphasis on the performance of the Gradient Boosting Regressor for crop yield prediction, the offline synchronization architecture, the RAG-based intelligent query system, the stakeholder-specific dashboards, and the legacy data import module. Qualitative evaluation of the system is supplemented by quantitative test set analysis and performance benchmarking.

\section{Machine Learning Yield Prediction Performance}

The optimized Gradient Boosting Regressor constitutes the core analytical engine of the system, trained to predict crop yields for major Ethiopian crops including teff, wheat, maize, barley, sorghum, sesame, chickpea, and coffee. The rigorous quantitative analysis of the model on an unseen test dataset demonstrated robust predictive performance.

\subsection{Quantitative Performance on the Test Set}

The most definitive assessment of the model's capabilities lies in its performance on the entirely independent test set, which was not used at any stage of model training or hyperparameter tuning. The quantitative metrics achieved by the Gradient Boosting model on this test set are detailed below.

\begin{table}[H]
\centering
\caption{Gradient Boosting Regressor Performance Metrics}
\label{tab:model-metrics}
\begin{tabular}{lc}
\toprule
\textbf{Metric} & \textbf{Value} \\
\midrule
Mean Absolute Error (MAE) & 863.54~kg/ha \\
Root Mean Square Error (RMSE) & 1539.96~kg/ha \\
Training Records & 1,700 \\
Test Records & 300 \\
Feature Dimension & 7 \\
Model Version & 20260520\_131732 \\
Model Type & GradientBoostingRegressor \\
Hyperparameters & n\_estimators=200, max\_depth=5, lr=0.1, subsample=0.8 \\
Training Data Source & CSV (sample\_yield.csv) \\
\bottomrule
\end{tabular}
\end{table}

The Mean Absolute Error of 863.54~kg/ha indicates that on average, the model's predictions deviate from the actual yield by approximately 864~kg/ha. This level of accuracy represents the typical magnitude of error a user can expect when using the system for yield prediction. For context, if a farmer expects a yield of 20,000~kg from a plot, the model's prediction would typically fall within approximately ±864~kg of the actual yield.

The Root Mean Square Error of 1539.96~kg/ha provides additional insight into the error distribution. The RMSE is higher than the MAE (approximately 1.78 times larger), which suggests that while most predictions are reasonably accurate, there are some predictions with larger errors. These larger errors likely correspond to outlier conditions such as extreme weather events or unusual agricultural practices that are underrepresented in the training data.

\subsection{Model Version Comparison}

The system's model versioning capability allows tracking of model performance across different training runs. Two model versions were produced during development:

\begin{table}[H]
\centering
\caption{Model Version Comparison}
\label{tab:model-versions}
\begin{tabular}{lcc}
\toprule
\textbf{Version} & \textbf{MAE (kg/ha)} & \textbf{RMSE (kg/ha)} \\
\midrule
20260512\_112525 & 2,424.24 & 3,408.94 \\
20260520\_131732 & 863.54 & 1,539.96 \\
\bottomrule
\end{tabular}
\end{table}

The significant improvement between version 1 (MAE 2,424.24) and version 2 (MAE 863.54) was achieved through improved data preprocessing, expanded training data, and optimized hyperparameter tuning. This represents a 64.4 percent reduction in MAE, demonstrating the importance of iterative model refinement.

\subsection{Feature Engineering Results}

The preprocessing pipeline successfully transformed raw agricultural product data into a format suitable for gradient boosting. The seven-feature vector (five numerical normalized, two categorical encoded) provided the model with sufficient information to capture the key factors influencing crop yields in Ethiopia.

The categorical encoding successfully captured eight crop types and two seasons. The label encoding approach was chosen over one-hot encoding because gradient boosting can effectively handle label-encoded categorical features through its tree-based splitting mechanism. The mapping of unseen categories to a default value of 0 ensures graceful degradation when the model encounters crop types or seasons not present in the training data.

The z-score normalization for numerical features ensured that all features contributed equally to the model's splitting decisions, preventing features with larger numerical ranges (such as area\_hectares with a range of 0--20) from dominating features with smaller ranges (such as temperature\_celsius with a range of 11--33).

\subsection{Qualitative Verification Through Web Application}

The trained model was integrated into the web application to verify its practical utility. Manual testing involved submitting prediction requests for each of the major crop types across different regions and seasons. Test inputs were chosen to incorporate variations in rainfall, temperature, and fertilizer amounts covering various conditions likely to occur in actual agricultural settings.

The system successfully processed all prediction requests and returned yield estimates with model version information. The predictions varied appropriately with input changes: increasing rainfall or fertilizer generally resulted in higher predicted yields, while extreme temperatures or low rainfall resulted in lower predictions. This behavior aligns with domain expectations for agricultural yield response to input variations.

The web interface displays the predicted yield prominently along with the model version used for the prediction. Users can view their prediction history and compare predictions across different input scenarios, enabling what-if analysis for agricultural planning.

\section{Offline Synchronization Performance}

The synchronization engine was tested by simulating data entry while the device was offline, followed by automatic synchronization when connectivity was restored. Test scenarios incorporated variations in network conditions, including complete loss of connectivity and intermittent connectivity with frequent drops.

\subsection{Synchronization Queue Management}

The sync queue system successfully tracked operations across five status levels:

\begin{table}[H]
\centering
\caption{Sync Queue Status Distribution}
\label{tab:sync-status}
\begin{tabular}{lc}
\toprule
\textbf{Status} & \textbf{Description} \\
\midrule
Pending & Operation queued, awaiting processing \\
Processing & Operation currently being processed \\
Completed & Operation successfully applied \\
Failed & Operation failed and requires retry \\
Conflict & Operation conflicted with server state \\
\bottomrule
\end{tabular}
\end{table}

The queue-and-check mechanism proved effective at managing synchronization during intermittent connectivity. The system correctly:
\begin{itemize}
  \item Stored operations with complete metadata including device ID, user ID, batch ID, entity type, operation type, and client timestamp
  \item Processed operations in batches via the batch endpoint
  \item Detected conflicts through timestamp comparison
  \item Allowed retry of failed operations with retry count tracking
  \item Updated status through the dedicated status update endpoint
\end{itemize}

\subsection{Device Management}

The system successfully manages field devices through a dedicated device registration and tracking system. Each device is identified by a unique hardware UUID and associated with a specific user. The device record tracks:
\begin{itemize}
  \item \textbf{Device UUID:} Hardware-identifying unique identifier
  \item \textbf{Device name:} Human-readable name for identification
  \item \textbf{User assignment:} Association with a specific field agent
  \item \textbf{Last sync timestamp:} When the device last synchronized data
  \item \textbf{Status:} Online, offline, or inactive
\end{itemize}

\section{RAG-Based Intelligent Query Performance}

The RAG query system was tested with natural language questions covering various aspects of agricultural product data including yield trends, regional comparisons, crop performance, and seasonal patterns.

\subsection{Query Processing Pipeline}

The RAG pipeline processes queries through three stages:

\subsubsection{Stage 1: Query Classification and Data Retrieval}

The system analyzes the user's question using keyword heuristics to determine which data sources to query. The keyword matching logic identifies query intent through the following patterns:

\begin{itemize}
  \item \textbf{Yield queries:} Keywords including yield, harvest, production, crop trigger retrieval from the imported\_records table with aggregation by crop name, region, and season.
  \item \textbf{Farmer queries:} Keywords including farmer, kebele, household trigger retrieval from the farmers table with distribution by region and woreda.
  \item \textbf{Observation queries:} Keywords including observation, field, health, pest trigger retrieval from the observations table joined with plots, farmers, and crop\_types.
  \item \textbf{Prediction queries:} Keywords including predict, forecast, model, ml trigger retrieval from the predictions table with model version grouping.
  \item \textbf{Trend queries:} Keywords including trend, decline, increase, change trigger yield data retrieval with temporal analysis.
  \item \textbf{Region queries:} Keywords including region, zone, woreda, or specific region names trigger region-specific data retrieval.
\end{itemize}

The system also supports explicit dataset filtering through the dataset\_type parameter, allowing users to restrict queries to imported records, field observations, farmer data, or all sources.

\subsubsection{Stage 2: Context Assembly}

Retrieved data is formatted into a structured context block with the following organization:
\begin{enumerate}
  \item System overview with summary statistics (total farmers, plots, observations, predictions)
  \item Historical yield data from imported records with crop/region/season aggregation
  \item Live field observations from the observations table with health and pest statistics
  \item Farmer distribution data by region and woreda
  \item ML prediction statistics with model version grouping
\end{enumerate}

Each section includes citations describing the data source and row count, enabling the LLM to reference specific data sources in its response.

\subsubsection{Stage 3: LLM Response Generation}

The assembled context is sent to the configured LLM along with a system prompt that provides:
\begin{itemize}
  \item Role definition: "agricultural product data analyst assistant for Ethiopia's Integrated agricultural product data System"
  \item Target audience: government officials, NGOs, researchers, traders
  \item Response guidelines: use only provided data, cite specific numbers, highlight patterns, acknowledge limitations
  \item Ethiopian context: main crops, seasons, regions
\end{itemize}

\subsection{LLM Provider Performance}

The system supports three LLM providers with automatic fallback:

\begin{enumerate}
  \item \textbf{Anthropic Claude (claude-haiku-4-5-20251001):} Primary provider optimized for speed and cost. Suitable for the moderate query volumes expected in agricultural applications.
  \item \textbf{Cerebras (llama3.1-70b):} Alternative provider with higher computational capacity, suitable for complex queries requiring deeper analysis.
  \item \textbf{OpenRouter (meta-llama/llama-3.1-8b-instruct):} Fallback provider for situations where both primary and alternative providers are unavailable.
\end{enumerate}

\subsection{System Prompt Design}

The RAG system uses a carefully designed system prompt that instructs the LLM to:
\begin{enumerate}
  \item Answer accurately using ONLY the data provided in the context
  \item Be specific by citing regions, crop names, and numerical values
  \item Highlight meaningful patterns such as declining yields or regional disparities
  \item Acknowledge when data is insufficient to answer the question
  \item Use plain language suitable for both technical and non-technical stakeholders
  \item Mention data limitations when relevant
\end{enumerate}

This prompt engineering ensures that responses are grounded in actual database records and provide actionable insights rather than generic agricultural advice.

\section{Stakeholder Dashboard Performance}

The stakeholder-specific dashboards implement role-based access control to provide appropriate data views for each user type. The dashboard controller dynamically filters data based on the caller's role, ensuring that each stakeholder group sees only the information relevant to their needs.

\subsection{Dashboard Data Visibility by Role}

\begin{table}[H]
\centering
\caption{Dashboard Payload Components by User Role}
\label{tab:payload-role}
\begin{tabularx}{\textwidth}{lX}
\toprule
\textbf{Role} & \textbf{Payload Components} \\
\midrule
Admin & summary, yieldSummary, farmerCounts, syncStatus, recentLogs, importStats, predictionStats \\
Supervisor & summary, yieldSummary, farmerCounts, syncStatus, recentLogs, importStats, predictionStats \\
Field Agent & summary, yieldSummary \\
Government & summary, yieldSummary, farmerCounts, syncStatus \\
NGO & summary, yieldSummary, farmerCounts, syncStatus, importStats, predictionStats \\
Trader & summary, yieldSummary \\
Researcher & summary, yieldSummary, farmerCounts, syncStatus, importStats, predictionStats \\
\bottomrule
\end{tabularx}
\end{table}

\subsection{Summary Statistics}

The dashboard summary provides a high-level overview of the system's data:

\begin{table}[H]
\centering
\caption{Dashboard Summary Statistics}
\label{tab:dashboard-summary}
\begin{tabularx}{\textwidth}{lX}
\toprule
\textbf{Statistic} & \textbf{Description} \\
\midrule
totalFarmers & Count of active farmer records (not soft-deleted) \\
totalPlots & Count of active plot records \\
totalExpectedYieldKg & Sum of expected yield from all observations \\
totalActualYieldKg & Sum of actual (harvested) yield from all observations \\
pendingSyncs & Count of sync queue entries still in pending status \\
\bottomrule
\end{tabularx}
\end{table}

\subsection{Frontend Visualization Components}

The dashboard frontend provides comprehensive data visualization through:

\begin{description}
  \item[AreaChart (Yield by Crop):] Displays yield distribution across crop types with area fill, enabling visual comparison of production volumes.
  \item[BarChart (Farmer Counts):] Shows farmer distribution across organizations or regions, facilitating resource allocation decisions.
  \item[PieChart (Observation Share):] Visualizes the proportion of observations by growth stage, health status, or crop type.
  \item[RadarChart (Yield Efficiency):] Compares yield performance across multiple dimensions for holistic assessment.
  \item[LineChart (Farmer Growth):] Tracks farmer registration trends over time for monitoring system adoption.
  \item[Interactive Map (Plot Locations):] Leaflet-based map displaying plot markers with popup information including farmer name, crop type, and area.
\end{description}

\subsection{Dashboard API Endpoints}

The dashboard controller provides specialized aggregation endpoints:

\begin{itemize}
  \item \texttt{GET /api/v1/dashboard}: Unified dashboard overview with role-filtered data
  \item \texttt{GET /api/v1/dashboard/farmer-counts}: Active farmer count per organization
  \item \texttt{GET /api/v1/dashboard/yield-summary}: Yield estimates aggregated by crop type
  \item \texttt{GET /api/v1/dashboard/sync-status}: Pending/failed/conflict sync counts per device
  \item \texttt{GET /api/v1/dashboard/health}: API health check with database connectivity status
\end{itemize}

\section{Legacy Data Import Performance}

The legacy data import module was tested with CSV files of varying sizes and complexities. The module successfully handles the complete import workflow from file upload to data persistence.

\subsection{Schema Detection Accuracy}

The automatic schema detection system matches column headers to 20 canonical database fields using over 80 known column aliases. The system successfully identified column mappings for a wide range of header variations including:

\begin{itemize}
  \item \textbf{Crop name:} crop, crop\_name, crop name, crop\_type, crop type, plant, item, commodity
  \item \textbf{Yield:} actual\_yield, actual yield, actual\_yield\_kg, yield, production, harvest, yield\_kg
  \item \textbf{Area:} area, area\_ha, hectares, area\_hectares, plot\_size, farm\_size
  \item \textbf{Rainfall:} rainfall, rainfall\_mm, precipitation, rain, annual\_rainfall, average\_rainfall\_mm\_per\_year
  \item \textbf{Temperature:} temperature, temp, temp\_celsius, temperature\_celsius, avg\_temp\_c, avg\_temp
  \item \textbf{Season:} season, growing\_season, period
  \item \textbf{Fertilizer:} fertilizer, fertilizer\_type, fertiliser, fertiliser\_amount, fertilizer\_kg
\end{itemize}

\subsection{Data Validation Rules}

The import module enforces comprehensive validation on each record:

\begin{longtable}{lll}
\caption{Import Validation Rules}
\label{tab:validation-rules} \\
\toprule
\textbf{Field} & \textbf{Rule} & \textbf{Error Message} \\
\midrule
\endfirsthead
\caption*{Table~\ref{tab:validation-rules}: Import Validation Rules (continued)} \\
\toprule
\textbf{Field} & \textbf{Rule} & \textbf{Error Message} \\
\midrule
\endhead
\midrule
\endfoot
\bottomrule
\endlastfoot
area\_hectares & 0 to 100,000 & "{value}" seems out of plausible range \\
actual\_yield\_kg & 0 to 1,000,000 & "{value}" seems out of plausible range \\
rainfall\_mm & 0 to 20,000 & "{value}" seems out of plausible range \\
temperature\_celsius & $-90$ to $60$ & "{value}" outside Earth surface range \\
year & 1900 to 2100 & "{value}" outside plausible range \\
All numeric fields & Must be numeric & "{value}" is not a valid number \\
Date fields & ISO or DD/MM/YYYY & "{value}" is not a recognisable date \\
\end{longtable}

\subsection{Import Processing Pipeline}

The complete import pipeline processes files through the following stages:
\begin{enumerate}
  \item \textbf{File parsing:} CSV files are parsed using the csv-parse library with streaming for memory efficiency. XLSX files are parsed using the ExcelJS library with worksheet extraction.
  \item \textbf{Schema detection:} Column headers are matched to canonical field names through the alias dictionary.
  \item \textbf{Record normalization:} Each row is normalized with string trimming, numeric coercion, date formatting, and season normalization.
  \item \textbf{Validation:} Each normalized record is checked against validation rules and flagged as valid, invalid, or skipped.
  \item \textbf{Report generation:} A validation report is produced with total, valid, invalid, and skipped counts plus per-row error details.
\end{enumerate}

\section{System Architecture Verification}

\subsection{API Completeness}

The system implements 108 RESTful API endpoints across 16 controllers, providing comprehensive coverage of all functional requirements. The endpoint distribution demonstrates the system's breadth:

\begin{table}[H]
\centering
\caption{API Endpoint Distribution}
\label{tab:api-distribution}
\begin{tabularx}{\textwidth}{lX}
\toprule
\textbf{Domain} & \textbf{Endpoints} \\
\midrule
Authentication & 12 (signup, login, tokens, password reset, email verification) \\
Profile & 3 (show, update, delete) \\
Dashboard & 5 (overview, farmer counts, yield, sync, health) \\
Organizations & 6 (CRUD + organization dashboard) \\
App Users & 8 (CRUD + activate/deactivate + devices) \\
Devices & 8 (CRUD + UUID lookup + sync marking) \\
Crop Types & 7 (CRUD + categories) \\
Farmers & 9 (CRUD + stats + plots + restore) \\
Plots & 8 (CRUD + nearby + restore) \\
Observations & 9 (CRUD + summary + harvest + restore) \\
Attachments & 6 (CRUD + by observation) \\
Sync Queue & 8 (CRUD + stats + batch + status + retry) \\
Audit Logs & 5 (CRUD + stats + by record) \\
Imports & 5 (upload + list + detail + schema + records) \\
Predictions & 4 (create + list + detail + ML health) \\
RAG & 2 (query + status) \\
Monetization & 8 (products + subscriptions + transactions + webhook) \\
Analytics & 2 (import trends + yield trends) \\
\bottomrule
\end{tabularx}
\end{table}

\subsection{Database Schema Verification}

The database schema was verified to contain 15 tables covering all required domains. The schema includes proper foreign key relationships, unique constraints, and indexes for query performance. Soft-delete support is implemented through deleted\_at timestamps for farmers, plots, and observations. Version-based optimistic locking is implemented for plots and observations through a version column that increments on each update.

\subsection{Middleware and Security Verification}

The system implements a multi-layered security architecture:

\begin{description}
  \item[Force JSON Response:] All requests are forced to return JSON responses for consistent API behavior.
  \item[Authentication:] The auth middleware blocks unauthenticated access to protected endpoints.
  \item[Email Verification:] Verified email middleware blocks access for users who have not verified their email.
  \item[Role Guard:] Role guard middleware checks user role against allowed roles for specific endpoints.
  \item[CORS:] Cross-Origin Resource Sharing middleware enables controlled frontend-backend communication.
  \item[CSRF Protection:] Shield middleware provides CSRF protection for state-changing requests.
\end{description}

\section{Discussion}

\subsection{Key Achievements}

The Gradient Boosting Regressor, achieving a test MAE of 863.54~kg/ha, provides a robust analytical core for crop yield prediction. The feature engineering approach using z-score normalization for numerical features and label encoding for categorical features proved effective for the agricultural product data set. The model versioning system enables continuous improvement as more training data becomes available.

The offline synchronization architecture provides a practical solution for agricultural product data management in low-connectivity environments. The server-side queue-and-check mechanism, combined with device tracking and timestamp-based conflict resolution, addresses the fundamental challenge of data collection in rural Ethiopia where internet connectivity is unreliable. The five-state status tracking provides transparency into the synchronization process.

The RAG-based intelligent query system represents a significant advancement in making agricultural product data accessible to non-technical users. By grounding LLM responses in actual database records through SQL-based retrieval and providing citations to specific data sources, the system achieves factual reliability that general-purpose language models cannot guarantee. The support for multiple LLM providers ensures operational resilience.

The stakeholder-specific dashboards with role-based access control successfully address the diverse information needs of different user groups. The dashboard controller's role-filtering mechanism ensures that government users can access macro-level statistics for policy planning, NGO users can identify priority intervention areas, traders can access market-relevant data, and researchers can access detailed records for analysis.

The legacy data import module reduces the barrier to digital adoption by enabling organizations to bulk-upload their existing spreadsheet records. The automatic schema detection system, with over 80 column header aliases, handles the variability in column naming conventions common in manually maintained agricultural records. The comprehensive validation system ensures data quality during import.

\subsection{Comparison with Related Work}

Compared to the offline-first data collection systems studied by Islam and colleagues \cite{islam2021}, the present system extends the queue-and-check pattern to agricultural-specific data with complex entity relationships. The synchronization engine handles four entity types (farmers, plots, observations, attachments) with full CRUD operations and maintains a comprehensive audit trail. The system also integrates with device management for tracking individual field devices.

Compared to existing RAG systems in agriculture, the present system differs by using SQL-based retrieval from structured databases rather than vector similarity search on unstructured documents. This approach is better suited for the structured nature of agricultural records where precise numerical aggregation and filtering are required. The system also supports multiple LLM providers for flexibility and cost optimization, unlike most RAG systems that are tied to a single provider.

\subsection{Limitations Encountered}

Several limitations were encountered during development and testing. The trained model achieved an MAE of 863.54~kg/ha, which while useful for regional planning, may not be sufficiently accurate for farm-level decision-making where more precise yield estimates are needed. The dataset of 2,000 records, while representative of major Ethiopian crops and regions, is relatively small for training highly accurate regression models. Expanding the dataset with more records across additional regions, years, and crop types would likely improve model performance.

The RAG system requires an active internet connection to communicate with LLM APIs, which limits its usefulness in completely offline scenarios. A future enhancement could include a locally-deployed smaller LLM for offline operation, though this would require additional computational resources.

The offline synchronization is currently implemented as a server-side queue rather than a full client-side Progressive Web App with IndexedDB storage. This means that the application still requires a browser session for data entry, albeit with deferred server synchronization. A true PWA implementation would enable data entry even when the browser tab is not actively connected to the server.

Image-based disease detection using convolutional neural networks, while mentioned in the project's planning phase, was not implemented in the current version. The observations module supports image attachments, providing the infrastructure for this enhancement, but the training and integration of an image classification model was beyond the scope of this initial implementation.

\subsection{Implications for Practice}

The Integrated agricultural product data System demonstrates the feasibility of combining modern web technologies, machine learning, and natural language processing to address agricultural product data management challenges in developing countries. The system's modular architecture and comprehensive API design provide a blueprint that can be adapted for other agricultural contexts beyond Ethiopia.

For agricultural extension services, the system provides a practical tool for digitizing their data collection workflows and accessing analytical insights. The offline synchronization capability ensures that the system remains useful even in areas with limited connectivity, addressing a key barrier to digital adoption in rural areas.

For policymakers, the system demonstrates how agricultural product data can be made accessible for evidence-based decision-making through role-based dashboards and natural language querying. The comprehensive audit logging ensures data accountability and supports regulatory compliance.
